\documentclass[11pt]{article}
\usepackage[localtheorems]{raeez-math-template}

\newtheorem{theorem}{Theorem}
\newtheorem{proposition}[theorem]{Proposition}
\newtheorem{lemma}[theorem]{Lemma}
\newtheorem{corollary}[theorem]{Corollary}
\theoremstyle{definition}
\newtheorem{definition}[theorem]{Definition}
\newtheorem{remark}[theorem]{Remark}

\providecommand{\C}{\mathbb{C}}
\providecommand{\Z}{\mathbb{Z}}
\providecommand{\Q}{\mathbb{Q}}
\providecommand{\cA}{\mathcal{A}}
\providecommand{\cC}{\mathcal{C}}
\providecommand{\cH}{\mathcal{H}}
\providecommand{\cM}{\mathcal{M}}
\providecommand{\cR}{\mathcal{R}}
\providecommand{\cT}{\mathcal{T}}
\providecommand{\fg}{\mathfrak{g}}
\providecommand{\Yplus}{Y^{+}}
\providecommand{\CoHA}{\mathrm{CoHA}}
\providecommand{\Stab}{\mathrm{Stab}}
\providecommand{\DT}{\mathrm{DT}}
\providecommand{\Li}{\mathrm{Li}}
\providecommand{\Ad}{\mathrm{Ad}}
\providecommand{\Aut}{\mathrm{Aut}}
\providecommand{\Rep}{\mathrm{Rep}}
\providecommand{\id}{\mathrm{id}}
\providecommand{\kcat}{\kappa_{\mathrm{cat}}}
\providecommand{\kch}{\kappa_{\mathrm{ch}}}
\providecommand{\kBKM}{\kappa_{\mathrm{BKM}}}
\providecommand{\kfib}{\kappa_{\mathrm{fiber}}}

\title{Kontsevich-Soibelman Wall-Crossing and the Positive Half
\texorpdfstring{$\Yplus$}{Y+}}
\author{Raeez Lorgat}
\date{}

\begin{document}
\maketitle

\begin{abstract}
Kontsevich-Soibelman wall-crossing is an ordered-product theorem in a
completed motivic quantum torus after Hall integration. A positive-half
CoHA presentation $\CoHA^+(X)\simeq\Yplus_X$ contributes to this theorem
through its completed Hall algebra, its integration character, and a
fixed shuffle realisation when such a realisation is known. Chamber
change preserves the ordered sector product after integration.
A chamber-to-chamber automorphism of $\Yplus_X$ is additional structure:
it must come from a Hall comparison, a Drinfeld double, a centre, or
stable-envelope modules. For a CY$_3$ input the specialised
Calabi-Yau-to-chiral output is $E_1$; the $E_2$ operation belongs to the
centre, double, or module layer. In particular
$\CoHA^+(\C^3)=\Yplus(\widehat{\mathfrak{gl}}_1)$ is a positive-half
statement. The full $\mathcal W_{1+\infty}$ object is reached only after
the double, central completion, and Fock or evaluation realisation have
been supplied.
\end{abstract}

\section{The Quantum-Torus Ambient}

\begin{definition}[Charge lattice and sector completion]
Let $\cC$ be a $3$-Calabi-Yau category equipped with a Bridgeland
stability condition $\sigma=(Z,\cA)$ on a connected open subset of
$\Stab(\cC)$. Put
\[
 \Gamma=K^{\mathrm{num}}_0(\cC),\qquad
 \langle \alpha,\beta\rangle
   =\chi(\alpha,\beta)
   =-\chi(\beta,\alpha).
\]
Equivalently, if one starts from a non-skew Euler pairing, this note
uses the half-skew convention
$\langle\alpha,\beta\rangle=(\chi(\alpha,\beta)-\chi(\beta,\alpha))/2$.
This is the normalization in which the elementary conifold pair has
$\langle\gamma_1,\gamma_2\rangle=1$.
For a strict convex sector $V\subset \C^*$ with boundary free of active
rays, set
\[
 \Gamma_V(\sigma)=\{\gamma\in\Gamma : Z_\sigma(\gamma)\in V\}.
\]
The completed quantum torus of the sector is
\[
 \widehat{\cT}_{\Gamma,V}
 =
 \widehat{\bigoplus}_{\gamma\in\Gamma_V(\sigma)}
 \Q(q^{1/2})\,\hat e_\gamma,
 \qquad
 \hat e_\alpha\hat e_\beta
 =
 q^{\langle\alpha,\beta\rangle/2}\hat e_{\alpha+\beta}.
\]
The completion is the monoid completion along effective charge in
$\Gamma_V(\sigma)$.
\end{definition}

\begin{definition}[Quantum dilogarithm factor]
Let
\[
 \Psi_q(x)=\prod_{r\geq0}(1-q^{r+1/2}x)^{-1}
\]
in the $x$-adic completion. If
$\bar\Omega(\gamma;q)=\sum_j a_{\gamma,j}q^{j/2}$ is the primitive
refined BPS invariant of charge $\gamma$ in a fixed orientation-data
normalization, define
\[
 U_\gamma
 =
 \prod_j \Psi_q(q^{j/2}\hat e_\gamma)^{a_{\gamma,j}}
 \in \widehat{\cT}_{\Gamma,V}^{\times}.
\]
The exponents $a_{\gamma,j}\in\Z$ are the integrality datum.
\end{definition}

\begin{lemma}[Integrality is the torus-automorphism condition]
\label{lem:integrality-automorphism}
Assume $\bar\Omega(\gamma;q)\in\Z[q^{\pm1/2}]$. For every
$\beta\in\Gamma$ the conjugation $\Ad_{U_\gamma}(\hat e_\beta)$ lies in
the rational localisation of the integral completed quantum torus
$\widehat{\cT}_{\Gamma,V}^{\Z}$. More explicitly, if
$a=\langle\gamma,\beta\rangle\geq0$, then
\[
 \Ad_{\Psi_q(q^{j/2}\hat e_\gamma)}(\hat e_\beta)
 =
 \hat e_\beta
 \prod_{r=0}^{a-1}
 (1-q^{j/2+r+1/2}\hat e_\gamma),
\]
and the case $a<0$ is obtained by inverting the same finite product.
\end{lemma}

\begin{proof}
The relation
$\hat e_\gamma\hat e_\beta=q^{\langle\gamma,\beta\rangle}
\hat e_\beta\hat e_\gamma$ gives
$f(\hat e_\gamma)\hat e_\beta
=\hat e_\beta f(q^{\langle\gamma,\beta\rangle}\hat e_\gamma)$ for every
formal series $f$. Taking $f(x)=\Psi_q(q^{j/2}x)$ gives
\[
 \Psi_q(q^{j/2+a}x)\Psi_q(q^{j/2}x)^{-1}
 =
 \prod_{r=0}^{a-1}(1-q^{j/2+r+1/2}x)
\]
for $a\geq0$. Multiplying over $j$ with integer exponents
$a_{\gamma,j}$ produces integer powers of the displayed factors. The
negative-pairing case follows by the same computation with the quotient
reversed.
\end{proof}

\section{Kontsevich-Soibelman Sector Products}

\begin{definition}[Primitive and rational BPS invariants]
The primitive refined invariants $\bar\Omega(\gamma;q)$ and the rational
DT invariants $\Omega(\gamma;q)$ are related, in the usual refined
multicover normalization, by
\[
 \Omega(\gamma;q)
 =
 \sum_{m\mid\gamma}
 \frac{1}{m}\,
 \frac{q^{1/2}-q^{-1/2}}{q^{m/2}-q^{-m/2}}\,
 \bar\Omega(\gamma/m;q^m).
\]
At $q=1$ this becomes
$\Omega(\gamma)=\sum_{m\mid\gamma}m^{-2}\bar\Omega(\gamma/m)$.
The ordered quantum-torus product is built from the primitive integral
factors; the rational invariants are logarithmic coefficients.
\end{definition}

\begin{theorem}[Kontsevich-Soibelman ordered product]
\label{thm:ks-sector-product}
Assume the support property and the motivic Hall integration map for the
sector $V$. For a stability condition $\sigma$ define
\[
 A_V(\sigma)
 =
 \prod_{\ell\subset V}^{\curvearrowright}
 \prod_{\gamma:\,Z_\sigma(\gamma)\in\ell}
 U_\gamma(\sigma),
\]
where rays are ordered by decreasing central-charge argument and the
factors on a fixed ray are grouped in the ray completion. If
$\sigma_0$ and $\sigma_1$ are connected by a path whose active rays stay
inside $V$, then
\[
 A_V(\sigma_0)=A_V(\sigma_1)
 \quad\text{in}\quad
 \widehat{\cT}_{\Gamma,V}^{\times}.
\]
The invariant is the ordered sector product in the completed torus.
\end{theorem}

\begin{proof}
In the motivic Hall algebra the Harder-Narasimhan factorisation of the
stack of objects in the sector $V$ is independent of the chamber after
the boundary of $V$ is fixed. Kontsevich-Soibelman integration sends
Hall convolution to the quantum-torus product, with the skew Euler form
producing the $q$-commutation factor. Applying the integration map to
the Harder-Narasimhan factorizations on the two sides gives the
displayed equality. The support property supplies the pronilpotent
convergence required for the sector completion.
\end{proof}

\begin{corollary}[Logarithmic form]
\label{cor:logarithmic-form}
Let $\fg_{\Gamma,V}$ be the completed quantum-torus Lie algebra under the
commutator. The logarithm
\[
 \Theta_V(\sigma)=\log A_V(\sigma)\in\widehat{\fg}_{\Gamma,V}
\]
is chamber-independent under the hypotheses of
Theorem~\ref{thm:ks-sector-product}. Different ray factorizations are
Baker-Campbell-Hausdorff presentations of the same element.
\end{corollary}

\section{The Positive-Half Comparison}

\begin{definition}[Positive-half CoHA datum]
For a toric CY$_3$ quiver with potential $(Q_X,W_X)$, a positive-half
CoHA datum consists of a heart, a positive cone $\Gamma^+$, orientation
data, and an associative critical Hall algebra
\[
 \cH_X^+
 =
 \bigoplus_{\gamma\in\Gamma^+}
 H^\bullet_{\mathrm{eq}}\bigl(\cM_\gamma,\phi_{\mathrm{Tr}W_X}\bigr),
 \qquad
 \cH_X^+\simeq \Yplus_X,
\]
together with the completed Hall integration character
\[
 I_\sigma:\widehat{\cH}_{X,V}^+\longrightarrow
 \widehat{\cT}_{\Gamma,V}.
\]
For $X=\C^3$, Schiffmann-Vasserot identify
$\cH_{\C^3}^+$ with the positive half
$\Yplus(\widehat{\mathfrak{gl}}_1)$.
\end{definition}

\begin{proposition}[Integrated character on \texorpdfstring{$\Yplus$}{Y+}]
\label{prop:yplus-integrated-character}
Let $X$ carry a positive-half CoHA datum on a sector $V$. For two
stability conditions $\sigma_0,\sigma_1$ as in
Theorem~\ref{thm:ks-sector-product}, the semistable Hall products
$\mathbf A_V(\sigma_i)\in\widehat{\cH}_{X,V}^+$ satisfy
\[
 I_{\sigma_0}\bigl(\mathbf A_V(\sigma_0)\bigr)
 =
 I_{\sigma_1}\bigl(\mathbf A_V(\sigma_1)\bigr)
 \quad\text{in}\quad
 \widehat{\cT}_{\Gamma,V}^{\times}.
\]
Thus the transported statement on $\Yplus_X$ is equality of integrated
characters. A genuine algebra automorphism of $\widehat{\Yplus}_X$ is an
additional lift $\widetilde A_V$ satisfying the intertwining condition
\[
 I_\sigma\circ\widetilde A_V
 =
 \Ad_{A_V(\sigma)}\circ I_\sigma .
\]
\end{proposition}

\begin{proof}
The first assertion is Theorem~\ref{thm:ks-sector-product} before and
after the identification $\cH_X^+\simeq\Yplus_X$. The integration map
is a Hall character, so it records the torus-valued DT series attached
to a completed Hall product. The last display is the defining square for
a lift from the torus to the positive-half algebra. KS wall-crossing
supplies the torus operator. Generators, relations, coproduct, Hopf
pairing, and a double for $\Yplus_X$ are additional Hall data.
\end{proof}

\begin{proposition}[Location of the \texorpdfstring{$E_2$}{E2} structure]
\label{prop:e2-location}
For CY$_3$ inputs the specialised output of $\Phi_3$ is $E_1$-chiral.
Braided, $R$-matrix, and $E_2$ structures used in the
Maulik-Okounkov framework live only after passing to an enhanced layer:
\[
 D(\Yplus_X),\qquad
 Z\bigl(\Rep(\Yplus_X)\bigr),\qquad
 \text{stable-envelope module categories},\qquad
 \text{central or Cartan completions}.
\]
Consequently
\[
 \CoHA^+(\C^3)=\Yplus(\widehat{\mathfrak{gl}}_1)
\]
is a positive-half statement. The vertex algebra
$\mathcal W_{1+\infty}$ belongs to the full-double, centre, completion,
and evaluation chain.
\end{proposition}

\begin{proof}
The CoHA product on a CY$_3$ quiver with potential is the associative
Hall product, hence the $E_1$ operation on the positive half. Braiding
appears after passing from a half algebra to representations, doubles,
centres, or stable-envelope correspondences. In the $\C^3$ model the
shuffle algebra gives the positive half of the affine Yangian; the
vertex algebra $\mathcal W_{1+\infty}$ is obtained after adding the
opposite half, completing the Cartan and central directions, and taking
Fock or evaluation realisations.
\end{proof}

\section{Rank-Two Wall: The Conifold Pentagon}

\begin{proposition}[The torus identity]
\label{prop:conifold-pentagon}
Let
$X_{\mathrm{con}}=\operatorname{Tot}(\mathcal O_{\mathbb P^1}(-1)
\oplus\mathcal O_{\mathbb P^1}(-1))$ and let
$\gamma_1,\gamma_2$ be the two simple charges in the elementary
rank-two wall with
$\langle\gamma_1,\gamma_2\rangle=1$. In the KS convention
\[
 \hat e_{\gamma_1}\hat e_{\gamma_2}
 =
 q^{1/2}\hat e_{\gamma_1+\gamma_2},
 \qquad
 \hat e_{\gamma_2}\hat e_{\gamma_1}
 =
 q^{-1/2}\hat e_{\gamma_1+\gamma_2},
\]
the elementary wall is the quantum pentagon
\[
 \Psi_q(\hat e_{\gamma_1})\,
 \Psi_q(\hat e_{\gamma_2})
 =
 \Psi_q(\hat e_{\gamma_2})\,
 \Psi_q(\hat e_{\gamma_1+\gamma_2})\,
 \Psi_q(\hat e_{\gamma_1}).
\]
The three exponents in this elementary $A_2$ scattering factor are
\[
 \varepsilon(\gamma_1)=
 \varepsilon(\gamma_2)=
 \varepsilon(\gamma_1+\gamma_2)=1.
\]
They are the exponents of this rank-two torus identity, separate from
the complete conifold CoHA primitive spectrum.
\end{proposition}

\begin{proof}
The active lattice at the elementary conifold wall is the $A_2$ lattice
with quantum variables $X=\hat e_{\gamma_1}$ and
$Y=\hat e_{\gamma_2}$ satisfying $XY=q^{1/2}\hat e_{\gamma_1+\gamma_2}$.
When the central-charge arguments of $\gamma_1$ and $\gamma_2$ exchange
order, consistency of the KS scattering diagram inserts the bound-state
ray $\gamma_1+\gamma_2$. Since
$q^{-1/2}XY=\hat e_{\gamma_1+\gamma_2}$, the standard compact form
$E(X)E(Y)=E(Y)E(q^{-1/2}XY)E(X)$ is exactly the displayed identity after
the change of variable between $E(z)=\prod_{r\ge0}(1+q^r z)$ and
$\Psi_q(x)=\prod_{r\ge0}(1-q^{r+1/2}x)^{-1}$. The sign convention for
the BPS index changes all three exponents coherently; the torus identity
is unchanged.
\end{proof}

\begin{remark}[Full conifold CoHA]
The Klebanov-Witten conifold critical CoHA has affine
$A_1^{(1)}$ real-root families $(n+1,n)$ and $(n,n+1)$ with primitive
supertrace $1$ in the parity-forgotten convention, and an imaginary
sector $(n,n)$ of multiplicity two before traceless projection. The
pentagon above isolates the first rank-two scattering wall. A full
conifold quantum-group construction requires the completed double,
Cartan part, Hopf pairing, and representation layer.
\end{remark}

\begin{corollary}[Positive-half reading]
On the conifold positive half the pentagon gives equality after Hall
integration. A lift to an automorphism of the conifold Yangian positive
half requires a chosen stable-envelope or double presentation and must
satisfy the intertwining square of
Proposition~\ref{prop:yplus-integrated-character}. The conifold
$\kch$ value remains
\[
 \kch(X_{\mathrm{con}})=1,
\]
the direct McKay value of the conifold chamber.
\end{corollary}

\section{Borcherds Constants and Yangian Branches}

\begin{proposition}[Separate invariant slots]
\label{prop:k3e-slots}
For $K3\times E$ the categorical, chiral, Borcherds, and fibre
invariants
occupy separate slots:
\[
 \kcat(K3\times E)=0,\qquad
 \kch(K3\times E)=0,\qquad
 \kch^{\mathrm{Heis}}(K3\times E)=3,
\]
\[
 \kBKM(\mathfrak{g}_{\Delta_5})=5,\qquad
 \kfib(K3)=24.
\]
For the primitive CHL series indexed by $N\in\{1,2,3,4,6\}$,
\[
 (c_N(0))=(10,8,6,4,2),\qquad
 \bigl(\kBKM(\Phi_N)\bigr)=(5,4,3,2,1),
\]
with $\kBKM(\Phi_N)=c_N(0)/2$ term by term.
\end{proposition}

\begin{proof}
The compact value $\kch(K3\times E)=0$ is the Hodge supertrace
$\sum_q(-1)^q h^{0,q}(K3\times E)$, and $\kcat(K3\times E)=0$ is the
same K\"unneth-multiplicative Euler characteristic of $\mathcal O$ on
the total space. The Heisenberg-Mukai value $3$ belongs to the chiral
specialisation. The BKM value is the Borcherds weight, computed by the
constant term $c_N(0)/2$. The fibre value $24$ is the Mukai-lattice
rank of the K3 fibre. KS wall-crossing supplies the chamber product in
the completed torus; these four values enter through the four named
constructions above.
\end{proof}

\begin{remark}[Two branches]
The Hall-Drinfeld/BKM branch starts from the Hall double of the
$K3\times E$ CoHA and its Borcherds denominator. The Mukai Yangian branch
starts from the self-mirror Mukai-lattice Yangian. The two constructions
may share lattice data, but their generators, doubles, and automorphic
outputs are distinct mathematical objects.
\end{remark}

\section{Data for a General \texorpdfstring{$Y_X^+$}{Y+X} Lift}

For a toric or reduced CY$_3$ target $X$, a statement stronger than
Proposition~\ref{prop:yplus-integrated-character} consists of the
following data.

\begin{enumerate}
\item A fixed heart, positive cone $\Gamma^+$, support property, and
sector completion.
\item A critical CoHA or shuffle presentation
$\cH_X^+\simeq\Yplus_X$ with orientation data.
\item A Hall integration character
$I_\sigma:\widehat{\cH}_{X,V}^+\to\widehat{\cT}_{\Gamma,V}$ compatible
with convolution and the quantum-torus product.
\item A lift of each sector factor through a Drinfeld double, centre,
stable-envelope module, or central completion, giving an automorphism
$\widetilde A_V$ of the completed target and satisfying
$I_\sigma\circ\widetilde A_V=\Ad_{A_V(\sigma)}\circ I_\sigma$.
\item An $E_2$ claim stated only on the lifted layer from item 4; the
positive half itself carries the $E_1$ Hall product.
\end{enumerate}

\section*{References}

\begin{thebibliography}{99}
\bibitem{KS2008}
M. Kontsevich and Y. Soibelman, \emph{Stability structures, motivic
Donaldson-Thomas invariants and cluster transformations},
arXiv:0811.2435.

\bibitem{JS2012}
D. Joyce and Y. Song, \emph{A theory of generalized Donaldson-Thomas
invariants}, Memoirs of the AMS 217 (2012).

\bibitem{MPS2011}
J. Manschot, B. Pioline and A. Sen, \emph{Wall crossing from Boltzmann
black hole halos}, Journal of High Energy Physics 2011.

\bibitem{SV2013}
O. Schiffmann and E. Vasserot, \emph{Cherednik algebras, W-algebras and
the equivariant cohomology of the moduli space of instantons on
$\mathbb A^2$}, Publications Math\'ematiques de l'IH\'ES 118 (2013),
arXiv:1202.2756.

\bibitem{NN2011}
K. Nagao and H. Nakajima, \emph{Counting invariant of perverse coherent
systems and its wall-crossing}, International Mathematics Research
Notices 2011.

\bibitem{DM2016}
B. Davison and S. Meinhardt, \emph{Cohomological Donaldson-Thomas
theory of a quiver with potential and quantum enveloping algebras},
Inventiones Mathematicae 221 (2020), arXiv:1601.02479.

\bibitem{MO2012}
D. Maulik and A. Okounkov, \emph{Quantum groups and quantum cohomology},
Ast\'erisque 408 (2019), arXiv:1211.1287.

\bibitem{Borcherds1998}
R. Borcherds, \emph{Automorphic forms with singularities on Grassmannians},
Inventiones Mathematicae 132 (1998).

\bibitem{GritsenkoNikulin1997}
V. Gritsenko and V. Nikulin, \emph{Siegel automorphic form corrections
of some Lorentzian Kac-Moody Lie algebras}, American Journal of
Mathematics 119 (1997), arXiv:alg-geom/9504006.

\bibitem{Gritsenko1999}
V. Gritsenko, \emph{Elliptic genus of Calabi-Yau manifolds and Jacobi
and Siegel modular forms}, Algebra i Analiz 11 (1999), number 5, 100-125;
St. Petersburg Mathematical Journal 11 (2000), number 5, 781-804,
arXiv:math/9906190.
\end{thebibliography}

\end{document}
